\begin{abstract}
Large language models produce em dashes at varying rates, and the observation that some models ``overuse'' them has become one of the most widely discussed markers of AI-generated text. Yet no mechanistic account of this pattern exists, and the parallel observation that LLMs default to markdown-formatted output has never been connected to it. We propose that the em dash is \textit{markdown leaking into prose}---the smallest surviving unit of the structural orientation that LLMs acquire from markdown-saturated training corpora. We present a five-step genealogy connecting training data composition, structural internalization, the dual-register status of the em dash, and post-training amplification. We test this with a two-condition suppression experiment across twelve models from five providers (Anthropic, OpenAI, Meta, Google, DeepSeek): when models are instructed to avoid markdown formatting, overt features (headers, bullets, bold) are eliminated or nearly eliminated, but em dashes persist---except in Meta's Llama models, which produce none at all. Em dash frequency and suppression resistance vary from 0.0 per 1,000 words (Llama) to 9.1 (GPT-4.1 under suppression), functioning as a signature of the specific fine-tuning procedure applied. A three-condition suppression gradient shows that even explicit em dash prohibition fails to eliminate the artifact in some models, and a base-vs-instruct comparison confirms that the latent tendency exists pre-RLHF. These findings connect two previously isolated online discourses and reframe em dash frequency as a diagnostic of fine-tuning methodology rather than a stylistic defect.
\end{abstract}
